\section{Limitations}
\label{sec:limits}

\begin{itemize}
  \item \textbf{Reduced setting; baselines that learned less.} T5-base, 10 epochs (100 in a few runs), free GPUs. With 10 epochs the
  published methods are at the level of guessing on two of five tasks right after training on them and
  learned less on 15 tasks; one step of their code was a substitute (Appendix~\ref{app:impl}). Nothing
  follows about these methods in their published setting.
  \item \textbf{Without instructions: a first result.} Two published methods; one sequence of five
  tasks with trained answers (three seeds), of development tasks; at 10 epochs most of the gap in the final score is
  learning, and the 100-epoch runs are one seed. The bound in advance is loose, and it grows with the number of later tasks. Inputs that occur
  in two tasks can at best be split. The audit holds per task and pair of times; the feature, the
  floor and $\nu$ were chosen on development data. Rewordings by people are untested.
  \item \textbf{Chosen tasks, few seeds.} The five tasks were chosen from those one published method
  forgot most; about 100 test inputs per task.
  \item \textbf{Zones are narrow under rewording and add the proof, not better routing}
  (Section~\ref{sec:many}): machine paraphrases are certified for their own task in
  \WidenRefTfivePLow{}\% to \WidenRefTfivePHigh{}\% of cases; none of five attempts to widen the zones met
  its bar; the best rule without zones routes about five points better than with zones in front. The reworded instructions are a
  machine-written development set.
  \item \textbf{Older tasks have priority.} Among \ManyTasks{} real tasks \ManyCollidingPct{}\%
  collided, mostly siblings, and \ManyAllHeld{} could not be learned under their instruction;
  \exactfirst{} was not trained at scale.
  \item \textbf{Certified inputs only.} Under medium and heavy rewording most inputs are uncertified and
  kept their answers without a guarantee; sixteen typos leave 45\% on their task.
  The zones need an instruction with a common opening.
  \item \textbf{Several tests are about routes only} (stress, edits, orders, other models, paraphrases,
  re-routing; one split and one order); word order is ignored; no backward transfer; one trained model
  family; memory grows with one adapter per task; exactness presupposes the same inference set-up
  (Appendix~\ref{app:method}).
\end{itemize}

\section{Conclusion}
\label{sec:conclusion}

With old adapters frozen, an answer changes only if its route changes, so forgetting is a question
about the rule that sets the route. Benchmarks hide this, because their
inputs begin with the instruction; without it the two published methods we ran changed
\AudSeedsGainloraChanged{} and \AudSeedsInfloraChanged{} of \AudSeedsCompared{} answers, a loss that was
not forced. Our main rule, a recogniser on frozen states, blocks no task and comes with an audit and a
loose bound in advance; plain adapters behind it beat the two published methods in this setting, by a
line fixed before the run; with their full epochs (one seed) those methods learned the tasks and then
forgot \AudFullGainloraFTInt{} and \AudFullInfloraFTInt{} points.
When inputs carry an instruction, the answer to an input inside a zone cannot change, for any model and
any later tasks; in a reduced setting the certified switch changed no answer on five and on 15 tasks
and learned as much as plain adapters. The zones are narrow under rewording, where the recogniser
reached \AudUThreeHeavyRidgeAPOne{} against \AudUThreeHeavyCertAPOne{} (three seeds).
All was measured with a small model and, in all runs but a few, a short budget.

